% H9-T25: full generic high-degree covers via a scheme of divisors and pencils.
% Run JAC1964-20261121; /root/h9_height; GPT-6 Astra Pro.

\subsection*{H9-T25. Full generic covers in many degrees above the special range}

\begin{theorem}[H9-T25]
Let $A/k$, $n=2m\ge4$, and its genuine Lefschetz pencil
with critical set $\Sigma$ be as in H9-T19. Suppose
\[
 g\ge1,\qquad N\ge2,\qquad N>2g-2,
 \qquad
 \bigl(2\mid N\ \hbox{or}\ (g\ge2\hbox{ and }3\mid N)\bigr).
                                                               \tag{H9.150}
\]
In every characteristic the full geometric generic
smooth connected genus-$g$, degree-$N$ cover of the
fixed $\mathbf P^1$, unramified above $\Sigma$, satisfies
\[
 h_L(z,w)=\frac1N h_K(\operatorname{cor}z,
                              \operatorname{cor}w),\qquad
 \operatorname{rad}h_L
       =\operatorname{cor}^{-1}(\operatorname{rad}h_K).
                                                               \tag{H9.151}
\]
Here $F$ is the algebraic closure of the function field
of the explicit irreducible parameter scheme below,
$K=F(t)$, and $L=F(C)$ for its generic cover. All
homologically trivial rational codimension-$m$ cycles
are admissible. For an abelian fourfold, accepting
032, the full homological height group is positive
semidefinite.

The full fixed-target smooth-cover stack in this
range is irreducible, and the parameter scheme
dominates it. The norm kernel is proved height-radical,
without a claim that it is zero in the Chow group.
No assertion is made for arbitrary special covers.
\end{theorem}

\begin{proof}
The geometric parameter argument works for every
$N\ge2$ with $N>2g-2$ whenever the required open
set is nonempty. The divisibility condition in
(H9.150) will provide the already checked family
whose whole new cohomology has no finite-orbit vector.

\paragraph{An irreducible scheme of all source curves.}
For $g\ge2$, let $B_g$ be the scheme of smooth
tricanonically embedded genus-$g$ curves. It is smooth
and irreducible over $k$ and has a universal smooth
projective curve $\mathscr C_g\to B_g$. These are the
actual Hilbert parameter and universal family of
Deligne--Mumford, \emph{The irreducibility of the space
of curves of given genus}, \S1, printed p.~78, with
irreducibility proved for every algebraically closed
field in \S3, printed pp.~92--96. Thus no
characteristic restriction from a tame Hurwitz-space
theorem is being introduced.

For $g=1$ instead use the smooth irreducible
Weierstrass parameter
\[
 B_1=\operatorname{Spec}
          k[a_1,a_2,a_3,a_4,a_6,\Delta^{-1}],
\]
with its universal elliptic curve and zero section.
Every pointed genus-one curve is represented locally
by this parameter in every characteristic: Stacks
Project, Lemma~105.6.1, tag \texttt{072T}.
An arbitrary smooth genus-one curve over an
algebraically closed field has a point, so forgetting
the chosen origin covers every curve needed here.
We use this pointed parameter, not an assertion
about an unpointed genus-one moduli scheme.

\paragraph{Divisors and a vector bundle of pencils.}
Let $D_{g,N}$ be the complement of the diagonals in
the $N$-fold fibre product of $\mathscr C_g$ over
$B_g$. It parametrizes a source curve with $N$
ordered distinct points. It is smooth and irreducible:
the map to $B_g$ is smooth with geometrically
irreducible nonempty fibres. Pull back the universal
curve to $D_{g,N}$ and call it
$p:\mathscr C\to D_{g,N}$. The marked sections give
the relative effective divisor
\[
                 D=P_1+\cdots+P_N,
                   \qquad \mathcal L=\mathcal O_{\mathscr C}(D).
                                                               \tag{H9.152}
\]
Thus there is an actual universal line bundle,
without any choice of a universal line on a coarse
Picard space. By Serre duality and $N>2g-2$,
\[
 H^1(C,\mathcal L|_C)=0,\qquad
 \dim H^0(C,\mathcal L|_C)=N+1-g
                                                               \tag{H9.153}
\]
on every geometric fibre. Coherent cohomology and
base change therefore makes $\mathcal E=p_*\mathcal L$
a vector bundle of rank $N+1-g$, with arbitrary
base change compatible with its sections.

Take the total space of $\mathcal E\oplus\mathcal E$.
It parametrizes two ordered sections $(s_0,s_1)$ of
$\mathcal L$. The condition that they have no common
zero is open, since the curve family is proper.
Such a pair defines a nonconstant morphism
\[
                        \pi=[s_0:s_1]:C\to\mathbf P^1,
                               \qquad \pi^*\mathcal O(1)=\mathcal L,
                                                               \tag{H9.154}
\]
of degree $N$. Constancy would require a nowhere
vanishing section of the positive-degree line
$\mathcal L$, and is impossible. The relative map
is finite locally free of degree $N$: it is proper
with finite fibres, and fibrewise flatness for the
maps of smooth curves gives flatness over the
relative target.

Let $U_{g,N}$ be the further open where this map
is unramified above $\Sigma$. This condition is
open by the support of relative differentials
over the finite set $\Sigma$. In particular the
maps there are separable, since $\Sigma$ is nonempty.
When nonempty, $U_{g,N}$ is a smooth irreducible
scheme with a universal cover of the required type.

It parametrizes every geometric cover in question.
Indeed, given a separable degree-$N$ map
$\pi:C\to\mathbf P^1$, choose a target value over
which it is unramified. Its inverse image is a
reduced divisor $D$ of degree $N$; order its points.
The line $\pi^*\mathcal O(1)$ is isomorphic to
$\mathcal O(D)$, and the two pullback coordinate
sections give (H9.154). Choose a tricanonical
framing when $g\ge2$, or an origin and Weierstrass
equation when $g=1$. This presents the cover by a
geometric point of $U_{g,N}$. Hence the full
fixed-target cover stack is the image of this
irreducible parameter, is irreducible, and is
dominated by it.

\paragraph{Putting a checked family inside this actual parameter.}
Under (H9.150), H9-T23a gives a smooth irreducible
parameter scheme $S$ and a family of degree-$N$
genus-$g$ covers, unramified over $\Sigma$, for
which the entire new middle-cohomology trace
kernel has no nonzero finite-orbit vector.
For even $N$ this uses the degree-two final step
and is valid for every $g\ge1$. For $3\mid N$
with $g\ge2$, it uses the degree-three final step.
These families work also when the characteristic
divides $N$.

We explain why the choices in the preceding
parameter construction preserve that monodromy
property. Choose a fixed target value which is
unramified for the geometric generic cover over
$S$, and shrink $S$ so that its inverse image is
finite \'etale of degree $N$. After taking a
connected component of the finite \'etale cover
which orders this fibre, we have marked sections
$P_1,\ldots,P_N$ and the relative divisor $D$.
The pullback of a coordinate section gives a
relative isomorphism
$\pi^*\mathcal O(1)\simeq\mathcal O(D)$.

For $g\ge2$, trivialize $p_*\omega^{\otimes3}$ on
a further nonempty open to obtain a tricanonical
framing. For $g=1$ use $P_1$ as the origin and
choose a Weierstrass equation on a nonempty open,
as permitted by tag \texttt{072T}. These choices
give an actual morphism
\[
                              S'\longrightarrow U_{g,N}.
                                                               \tag{H9.155}
\]
The new-kernel sheaf on $S'$ still has no nonzero
finite-orbit vector. Finite-index restriction
preserves this property: a finite orbit under
a subgroup of finite index gives a finite orbit
under the entire group by taking finitely many
cosets. Restriction to a nonempty open of a normal
connected scheme also preserves it, since the
induced fundamental-group map is surjective.
This accounts for both the finite \'etale choices
and the open shrinkings. In particular (H9.155)
proves $U_{g,N}$ is nonempty.

\paragraph{The universal projector and the height identity.}
Over $U_{g,N}$ take
\[
 \mathscr X=(\widetilde A\times U_{g,N})
                  \times_{\mathbf P^1\times U_{g,N}}\mathscr C,
 \qquad Q:\mathscr X\to\widetilde A\times U_{g,N}.
                                                               \tag{H9.156}
\]
It is a smooth proper family of $n$-folds, since
the covers are unramified over all critical pencil
values. The map $Q$ is finite flat of degree $N$.
Consequently
\[
                           1-\frac1N Q^*Q_*
                                                               \tag{H9.157}
\]
cuts out a lisse trace-kernel summand of its
middle cohomology. Pulling it back along (H9.155)
gives the complete checked trace kernel on $S'$.
If a vector had finite orbit on $U_{g,N}$, its
unchanged fibre vector would have finite orbit
on $S'$. Thus the universal trace kernel has
no nonzero finite-orbit vector.

At $F=\overline{k(U_{g,N})}$ every algebraic
middle class has finite monodromy orbit. It
follows that for every codimension-$m$ rational
Chow cycle $Z$ on the geometric generic model,
\[
                    \left[Z-\frac1NQ^*Q_*Z\right]=0.
                                                               \tag{H9.158}
\]
This uses the actual finite-flat map of smooth
models and does not assume that algebraic cycles
descend in Chow under a generic parameter
extension. Admissible pushforward and pullback,
the intersection projection formula, and (H9.158)
give (H9.151), exactly as in (H9.149).
Corestriction is surjective over $\mathbf Q$
because $\operatorname{cor}\operatorname{res}=N$;
this proves the exact radical. Admissibility is
H9-T12(iii), and the fourfold sign follows from
H9-T8 under 032.
\end{proof}

\paragraph{Scope and remaining degree restrictions.}
This theorem applies, for example, to every even
degree $N>2g-2$ in every genus $g\ge1$, and to every
degree divisible by three above that bound when
$g\ge2$. It promotes the previously constructed
composite families to the full generic cover,
using a scheme which parametrizes every cover in
the degree range. The proof does not settle odd
prime degrees at least five, degree three in
genus one, or the general special-line-bundle
range $N\le2g-2$ beyond the cases already handled
by T20, T22, and T24.

\paragraph{Sources and status.}
Deligne--Mumford, \emph{Publications math\'ematiques
de l'IH\'ES} 36 (1969), pp.~75--109, was directly
read at the Hilbert parameter construction p.~78
and Corollary~1.7, p.~83, proving smoothness,
and the full irreducibility proof \S3,
pp.~92--96; the strengthened fine-level result
is Theorem~5.15, pp.~108--109:
\url{https://www.numdam.org/item/PMIHES_1969__36__75_0.pdf}.
The all-characteristic Weierstrass presentation
was directly read in Stacks Project
Lemma~105.6.1,
\url{https://stacks.math.columbia.edu/tag/072T}.
The vector-bundle consequence of (H9.153) also follows
directly from EGA~III$_2$, Theorem~7.7.5(I--II),
printed pp.~67--68 (continuous pp.~199--200), and
Proposition~7.8.4, printed pp.~73--74 (continuous
pp.~205--206). The bases here are smooth and reduced;
$H^1=0$ and Riemann--Roch give constant $h^0$, so
Proposition~7.8.4(e)$\Rightarrow$(d) gives local
freeness and tensor base change, while
Theorem~7.7.5(II)(d) gives arbitrary scheme base change:
\url{https://www.numdam.org/item/PMIHES_1963__17__5_0.pdf}.
The ordered-divisor parameter,
its dominance, and the monodromy-to-height
application are deductions proved here.
